% ===== SECTION 2: RELATED WORK =====

\section{Related Work}

\paragraph{Agent benchmarks.}
The rapid development of tool-using LLMs has produced numerous benchmarks. ToolLLM \cite{qin2023toollm} provides 16,000+ real-world APIs for evaluating tool selection and orchestration. API-Bank \cite{li2023apibank} evaluates API call accuracy across 73 APIs. AgentBench \cite{liu2024agentbench} spans eight environments including operating systems and web browsing. WebArena \cite{zhou2024webarena} evaluates autonomous agents in realistic, self-hosted web environments. SWE-bench \cite{jimenez2024swebench} targets software engineering with real GitHub issues. $\tau$-bench \cite{yao2024tau} introduces the pass$^k$ metric for measuring multi-trial consistency, showing that even state-of-the-art agents succeed on fewer than half the tasks at pass@1 and are highly inconsistent (pass$^8 < 25\%$ in retail).

However, these benchmarks share a common limitation: they report single-run success rates under idealized conditions, ignoring the production realities of varying inputs, infrastructure failures, and safety requirements. Table~\ref{tab:benchmark_comparison} summarizes the scale and evaluation protocol of the major agent benchmarks alongside our suite. Two structural contrasts stand out. First, scale: the community-standard benchmarks evaluate hundreds to thousands of tasks (WebArena: 812; $\tau$-bench: 165; SWE-bench: 2{,}294), whereas we trade task count for multi-dimensional measurement---each of our nine models is additionally evaluated on repeated runs, perturbed inputs, injected faults, and adversarial safety prompts, none of which the capability-scale benchmarks administer. Second, model coverage: every benchmark above evaluates frontier or large open-weight models, and $\tau$-bench explicitly excludes models under 13\,B as too weak for its difficulty. Our suite is the only one of the group that evaluates a multi-model cohort of sub-10\,B models on reliability dimensions at all.

\begin{table*}[t]
\centering
\caption{Comparison of our evaluation with major agent benchmarks. ``Reliability'' lists the reliability-oriented dimensions each work measures; ``Sub-10\,B'' indicates whether models under 10\,B parameters are included in the evaluated cohort.}
\label{tab:benchmark_comparison}
\small
\setlength{\tabcolsep}{4pt}
\begin{tabular}{lllllc}
\toprule
\textbf{Benchmark} & \textbf{Domain} & \textbf{Scale} & \textbf{Models} & \textbf{Reliability} & \textbf{Sub-10\,B} \\
\midrule
ToolLLM \cite{qin2023toollm} & 16{,}464 RapidAPI tools & 16k+ APIs & 4 & None & No \\
API-Bank \cite{li2023apibank} & 73 APIs & 753 dialogues & 4 & None & No \\
AgentBench \cite{liu2024agentbench} & 8 environments & 1{,}091 test queries & 27 & None & Few \\
WebArena \cite{zhou2024webarena} & 4 self-hosted websites & 812 tasks & 6 agents & None & No \\
SWE-bench \cite{jimenez2024swebench} & Real GitHub repos & 2{,}294 issues & 20+ & None & No \\
$\tau$-bench \cite{yao2024tau} & Retail, airline & 165 tasks & 7 & pass$^k$, cost & No \\
ReliabilityBench \cite{gupta2026reliabilitybench} & 4 domains & 1{,}280 episodes & 2 & $R(k,\varepsilon,\lambda)$ & No \\
Science of Reliability \cite{rabanser2026science} & 2 benchmarks & -- & 15 & 12 metrics, 4 dims & No \\
\midrule
\textbf{This work} & \textbf{8 categories} & \textbf{31 + 14 tasks} & \textbf{9} & \textbf{4 dims + composite} & \textbf{Yes} \\
\bottomrule
\end{tabular}
\end{table*}

\paragraph{Reliability evaluation.}
Several recent works have begun treating reliability as a first-class evaluation dimension. ReliabilityBench \cite{gupta2026reliabilitybench} introduces a three-dimensional reliability surface $R(k, \varepsilon, \lambda)$ capturing consistency, robustness, and fault tolerance, and evaluates it on two large models. The Science of AI Agent Reliability \cite{rabanser2026science} proposes twelve metrics across four dimensions---consistency, robustness, predictability, and safety---and evaluates 15 frontier models across two benchmarks. Claw-Eval \cite{claw2026eval} adds trajectory-aware grading to evaluate completion, safety, and robustness across 14 frontier models. A complementary diagnostic framework by Huang et al. \cite{huang2026diagnostic} proposes a 12-category error taxonomy for tool invocation failures and evaluates both open-weight and proprietary models, finding that procedural reliability---particularly tool initialization---is the primary bottleneck for smaller models.

Two further lines of work inform our methodology. Khanal et al. \cite{khanal2026longhorizon} introduce a reliability-science framework for long-horizon agents, showing that capability and reliability rankings diverge substantially as task duration increases---reinforcing our finding that accuracy alone is an insufficient deployment signal. And AgentProp-Bench \cite{gurram2026agentprop} validates LLM judges against human labels, finding that rejection and recovery from tool-use errors are statistically independent capabilities---consistent with our fault tolerance results, where schema-drift recovery is the most failure-prone mode and is achieved reliably only by the code-specialized model.

Crucially, every one of these studies evaluates only large frontier models (GPT-4o, GPT-5, Claude Opus, Gemini Pro, etc.). None include models under 10 billion parameters.

\paragraph{Safety evaluation.}
Safety-oriented red-teaming has likewise focused on frontier models. Universal and transferable adversarial attacks \cite{zou2023adversarial}, black-box jailbreaking via iterative query refinement \cite{chao2024jailbreaking}, and the standardized HarmBench framework \cite{mazeika2024harmbench} establish the threat model and evaluation methodology for refusal behavior at scale. Whether open-weight small models exhibit even basic refusal and scope-preservation behavior under agentic tool-use prompting remains unmeasured---a gap our safety dimension addresses.

\paragraph{Small model evaluation.}
The evaluation of small models has focused on capability rather than reliability. AgentFloor \cite{karmakar2026agentfloor} tests 16 open-weight models from 0.27\,B to 32\,B on a six-tier capability ladder, finding that small models match GPT-5 on structured tasks but fall short on long-horizon planning. The ``Right for Wrong Reasons'' study \cite{right2026wrong} reveals that 50--69\% of correct answers from 7--9\,B models contain fundamentally flawed reasoning, but focuses on static reasoning traces rather than multi-turn tool use.

TinyLLM \cite{tinyllm2026edge} evaluates sub-3\,B models on function-calling benchmarks, finding that medium-scale models (1--3\,B) achieve up to 65.74\% accuracy. EdgeVox \cite{edgevox2026} evaluates 18 SLMs for tool calling on edge devices. The ``Harness Design'' study \cite{cho2026harness} shows that with proper scaffolding, 2--3\,B models can achieve 95.2\% task success rates on a narrow set of 24 tasks. Lee et al. \cite{lee2025patool} identify \emph{schema misalignment}---where models hallucinate tool names absent from provided schemas---as a key failure mode, and propose adapting tool schemas to match models' pretraining distributions to improve accuracy by up to 17\%.

These studies provide valuable insights into what small models \emph{can do}, but not \emph{how reliably} they do it.

\paragraph{Our position.}
Table~\ref{tab:benchmark_comparison} situates our evaluation against the major agent benchmarks, and Table~\ref{tab:position} positions us in the reliability literature. To our knowledge (as of August 2026), this is the first study that combines (a) multiple small models (\textless10\,B), (b) a comprehensive reliability framework spanning multiple dimensions, and (c) tool-use agent tasks with multi-step execution.

\begin{table*}[t]
\centering
\caption{Positioning of our work relative to existing literature. 
         ``Reliability dims.'' indicates how many reliability dimensions are evaluated.}
\label{tab:position}
\small
\begin{tabular}{lccc}
\toprule
\textbf{Study} & \textbf{Models} & \textbf{Reliability dims.} & \textbf{Small models?} \\
\midrule
ReliabilityBench \cite{gupta2026reliabilitybench} & 2 & 3 & No \\
Science of Reliability \cite{rabanser2026science} & 15 & 4 & No \\
Claw-Eval \cite{claw2026eval} & 14 & 3 & No \\
AgentFloor \cite{karmakar2026agentfloor} & 16 & 0 (capability only) & Yes \\
Right-for-Wrong \cite{right2026wrong} & 3 & 1 (reasoning) & Yes \\
TinyLLM \cite{tinyllm2026edge} & sub-3\,B (4 families) & 0 (capability only) & Yes \\
\midrule
\textbf{This work} & \textbf{9} & \textbf{4} & \textbf{Yes} \\
\bottomrule
\end{tabular}
\end{table*}
